% Stage 8 restored fragment — 2.4.1 Sommerfeld radiation condition (archive lines 4939-5230)
\subsection{Sommerfeld radiation condition}
\label{sec:ch241}
Before Sommerfeld asymptotics are written, the physical ambiguity must be named. Harmonic
Helmholtz solutions on open domains superpose $e^{-\jj k_0 r}/r$ outgoing and
$e^{+\jj k_0 r}/r$ incoming spherical waves with equal standing rights in the volume
equations alone \cite{sommerfeld1949}. A near-field scan or pattern cut on
$\Gamma_R$ cannot distinguish these completions without an explicit outgoing selection: wattmeters
and probes then read mixtures of export, storage, and energy arriving from the computational or
chamber boundary \cite{hansen1988}. Sommerfeld radiation condition is the asymptotic
statement that restores laboratory identification---radiation launched locally and leaving the
domain---before any uniqueness argument or far-zone integral is evaluated
\cite{sommerfeld1949}.

Before any asymptotic formula is written, the range engineer already decides whether the
measurement enclosure behaves as an open exterior domain \cite{hansen1988}.
A compact range places the device under test inside an anechoic chamber whose absorbing
linings are designed to let export escape instead of reflecting back as multipath
\cite{collin2001}. Pattern probes on a large sphere $\Gamma_R$ and shell wattmeters report
tangential field and power only when chamber absorption and computational termination both
approximate free-space outgoing behavior \cite{poynting1884}. Incomplete
absorption produces ripple on azimuth cuts; unstable spherical near-field transforms when
spectra are truncated show the same branch failure at the data-processing stage
\cite{hansen1988}. Until outgoing selection is restored, $\Gamma_R$ meters read cavity
storage and standing-wave content rather than watt-bearing export \cite{maxwell1873}.
Sommerfeld radiation condition below is the asymptotic mathematical statement of what the
range engineer already assumes when attributing a measured pattern to the device under test
rather than to energy arriving from infinity \cite{sommerfeld1949}.
An open antenna measurement is meaningful only if the field on the range can be identified
with radiation from the device under test rather than with energy arriving from infinity.
Sommerfeld outgoing selection enforces that identification: disturbances launched locally
propagate outward and are not returned from the far boundary of the model
\cite{sommerfeld1949}. When the condition is omitted, the same near-field scan or pattern
cut admits formal completions with incoming or standing-wave content, and inversion may
assign equivalent sources that would require excitation mechanisms absent from the physical
problem and therefore lie outside the admissible solution class \cite{devaney2004}.
Free-space kernels
Eqs.~\eqref{eq:ch2.free-space-scalar-green} and~\eqref{eq:ch2.free-space-dyadic-green}
supply the asymptotic templates; outer-boundary classes from Chapter~\ref{ch:01} are
presupposed. Uniqueness proofs and reconstruction limits belong to
Subsections~\ref{sec:ch243} and~\ref{sec:ch244}.
Let \(\Omega_{\mathrm{ext}}\) denote an exterior region surrounding compact \(\operatorname{supp}\Jimp\subset\Omega_s\), with \(\omega>0\) fixed and vacuum
constants \(\varepsilon_0,\mu_0\) in the far zone. Spherical coordinates
\((r,\theta,\phi)\) are centered so that \(\Omega_s\subset B_R(0)\) for some
\(R_0>0\). Fields are phasors under \(\exp(-\jj\omega t)\); outgoing means
\(e^{-\jj k_0 r}/r\) decay with \(k_0=\omega\sqrt{\mu_0\varepsilon_0}\) from
Eq.~\eqref{eq:ch2.helmholtz-symbol} \cite{stratton1941}. Interior interfaces and
finite domains use Eq.~\eqref{eq:ch1.outer-boundary-classes}; Sommerfeld completion
applies on artificial or physical spheres at large \(r\) \cite{harrington2001}.
On $\Gamma_R$ with $R\gg\lambda_0$, a range engineer expects the measured phase to advance as
$k_0 r$ increases---outgoing fronts receding from the enclosure---and the amplitude to fall no
faster than $r^{-1}$ in three dimensions \cite{jackson1999}. That observation is what
Sommerfeld selection encodes at infinity: the scalar template states what every outgoing Helmholtz
response must look like at large distance from a localized source. Phase fronts recede with the
observer; amplitude decay compatible with energy export to infinity is $r^{-1}$ in three
dimensions \cite{jackson1999}. For scalar potentials and Helmholtz kernels, outgoing radiation
requires, with $r=|\mathbf{r}|$ the observation radius from an origin enclosing the source,
$k_0=\omega\sqrt{\mu_0\varepsilon_0}$ the vacuum wavenumber (rad\,m$^{-1}$), and
$A(\theta,\phi,\mathbf{r}')$ the angular amplitude that carries directional information while
$\mathbf{r}'$ labels the fixed source point during the limit $r\to\infty$:
\begin{equation}
\label{eq:ch2.sommerfeld-scalar-asymptotic}
G(\mathbf{r},\mathbf{r}')
=
A(\theta,\phi,\mathbf{r}')\,\frac{e^{-\jj k_0 r}}{r}
\;+\;
\mathcal{O}(r^{-2}),
\qquad
r=|\mathbf{r}|\to\infty,
\end{equation}
The remainder \(\mathcal{O}(r^{-2})\) collects subleading reactive and curvature corrections
that need not vanish at finite range \cite{stratton1941}. Equation
\eqref{eq:ch2.sommerfeld-scalar-asymptotic} is the Chapter~\ref{ch:02} scalar
Sommerfeld condition: it excludes incoming \(e^{+\jj k_0 r}/r\) and standing
\(e^{\pm\jj k_0 r}/r\) superpositions that degenerate exterior uniqueness
\cite{harrington2001}. The free-space template Eq.~\eqref{eq:ch2.free-space-scalar-green}
satisfies Eq.~\eqref{eq:ch2.sommerfeld-scalar-asymptotic} for fixed \(\mathbf{r}'\)
\cite{stratton1941}. Physically, Eq.~\eqref{eq:ch2.sommerfeld-scalar-asymptotic}
states that monopole-like responses propagate outward with phase fronts
\(k_0 r=\mathrm{const}\) and amplitude decay \(r^{-1}\), the only radial
combination compatible with energy flux to infinity in three dimensions
\cite{jackson1999}.
Vector fields must satisfy a transverse impedance relation at infinity: the electric and
magnetic components that carry export must stand in the same ratio as a plane wave
propagating along \(\hat{\mathbf{r}}\), with vacuum impedance \(\eta_0\) as the proportionality
constant \cite{stratton1941}. For harmonic Maxwell fields in vacuum far from sources,
\begin{equation}
\label{eq:ch2.sommerfeld-field-condition}
\mathbf{r}\times\E
\;+\;
\eta_0\,(\mathbf{r}\times\H)\times\hat{\mathbf{r}}
=
\mathbf{o}(1),
\qquad
r\to\infty,
\end{equation}
with \(\eta_0=\sqrt{\mu_0/\varepsilon_0}\) the vacuum wave impedance,
\(\hat{\mathbf{r}}=\mathbf{r}/r\) the unit radial direction, and \(\mathbf{o}(1)\) denoting
any vector whose Cartesian components vanish in the limit \(r\to\infty\)
\cite{stratton1941}. Physically, Eq.~\eqref{eq:ch2.sommerfeld-field-condition} requires that
the transverse electric field and the transverse magnetic field on a distant sphere lock
phase and magnitude as if energy were leaving the domain along \(\hat{\mathbf{r}}\), rather
than arriving from infinity \cite{harrington2001}. Equation~\eqref{eq:ch2.sommerfeld-field-condition}
is the Chapter~\ref{ch:02} vector Sommerfeld radiation condition. It implies
transverse radiative fields with asymptotic ratio
\begin{equation}
\label{eq:ch2.sommerfeld-impedance-ratio}
\H_{\mathrm{rad}}
\sim
\eta_0^{-1}\,\hat{\mathbf{r}}\times\E_{\mathrm{rad}},
\qquad
r\to\infty,
\end{equation}
and matches the outgoing-trace condition Eq.~\eqref{eq:ch1.radiation-trace-condition}
from Subsection~\ref{sec:ch01.1.2} when tangential traces are taken on a sphere of
large radius \cite{stratton1941}. Mathematically,
Eq.~\eqref{eq:ch2.sommerfeld-field-condition} selects the subspace on which the
exterior Laplacian/Helmholtz resolvent coincides with retarded Green functions
Eq.~\eqref{eq:ch2.retarded-harmonic-kernel-theorem}; physically, it enforces
power flow outward rather than inward at infinity \cite{harrington2001}.
Far-zone transversality and propagation. Sommerfeld-admissible fields decompose at leading order as
\begin{equation}
\label{eq:ch2.far-zone-transverse-radiation}
\E(\mathbf{r})
=
\E_0(\theta,\phi)\,\frac{e^{-\jj k_0 r}}{r}
\;+\;\mathcal{O}(r^{-2}),
\qquad
\E_0\cdot\hat{\mathbf{r}}=0,
\end{equation}
with \(\E_0\) the radiation amplitude pattern and \(\E_0\cdot\hat{\mathbf{r}}=0\)
the transversality constraint \cite{stratton1941}. Equation
\eqref{eq:ch2.far-zone-transverse-radiation} connects Sommerfeld completion to
\(\Pi_{\mathrm{prop}}\) from Subsection~\ref{sec:ch224}: far-zone content lies on
the propagating sheet with real \(k_r=k_0\) and outgoing phase
\cite{harrington2001}. Evanescent and reactive components may still be large at
finite \(r\) while satisfying Eq.~\eqref{eq:ch2.sommerfeld-field-condition}; regime
qualification Eq.~\eqref{eq:ch2.regime-operator-screen} remains separate from
Sommerfeld selection \cite{stratton1941}.
\begin{figure}[t]
\centering
\ChOneFigBlock{%
\begin{tikzpicture}[ch1fig, node distance=5mm and 7mm]
 \node[ch1box] (src) {Compact\\$\Omega_s$};
 \node[ch1box, right=16mm of src] (near) {Near/reactive\\$kr\lesssim 1$};
 \node[ch1box, right=16mm of near] (far) {Sommerfeld\\$e^{-jk_0 r}/r$};
 \node[ch1box, below=5mm of near] (ev) {$\Pi_{\mathrm{ev}}$};
 \draw[ch1flow] (src) -- (near) -- (far);
 \draw[ch1flow] (near.south) -- (near.south |- ev.north);
 \draw[ch1flow] (near.south |- ev.north) -- (ev.north);
\end{tikzpicture}%
}
\caption{Sommerfeld far zone Eq.~\eqref{eq:ch2.sommerfeld-field-condition}: compact
sources drive near fields with evanescent content; at large \(r\) only outgoing
transverse \(e^{-jk_0 r}/r\) waves satisfy exterior completion.}
\label{fig:ch2.sommerfeld-far-zone}
\end{figure}
Equation~\eqref{eq:ch2.sommerfeld-field-condition} is a local impedance statement at
infinity: transverse \(\E\) and \(\H\) must satisfy the plane-wave ratio
Eq.~\eqref{eq:ch2.sommerfeld-impedance-ratio} with \(\eta_0\) from
Eq.~\eqref{eq:ch1.vacuum-impedance-wavenumber}. Mathematically, it selects the
eigensheet of the exterior Helmholtz resolvent with outgoing radial phase; physically,
it enforces net power flow toward infinity rather than inward from infinity. On a
truncation sphere \(\Gamma_R\) with \(R=8\lambda_0\) enclosing a source with \(k_0 a=1.5\),
Silver--M\"uller surrogate Eq.~\eqref{eq:ch2.silver-muller-trace-approx} typically
introduces spurious reflected power below \(1\%\) of \(P_{\mathrm{rad}}\) while reactive
near fields at \(r=1.2a\) may differ by \(\mathcal{O}(1)\) from the exact Sommerfeld solution
\cite{harrington2001}. Sommerfeld completion is therefore necessary but not
sufficient for transport claims---regime qualification Eq.~\eqref{eq:ch2.regime-operator-screen}
remains mandatory \cite{stratton1941}. Figure~\ref{fig:ch2.sommerfeld-far-zone} separates Sommerfeld completion from regime qualification. Subsection~\ref{sec:ch241} fixes the asymptotic \emph{boundary at
infinity}; Subsection~\ref{sec:ch224} and Section~\ref{sec:ch01.2.4} govern when
far-zone samples may be interpreted as transport \cite{stratton1941}. Green-function asymptotics. Dyadic kernels from Subsection~\ref{sec:ch231} inherit scalar Sommerfeld decay.
For \(r=|\mathbf{r}-\mathbf{r}'|\to\infty\) at fixed \(\mathbf{r}'\),
\begin{equation}
\label{eq:ch2.sommerfeld-dyadic-asymptotic}
\overline{\mathbf{G}}_{E,0}(\mathbf{r},\mathbf{r}')
=
\overline{\mathbf{G}}_{\infty}(\theta,\phi,\mathbf{r}')\,
\frac{e^{-\jj k_0 r}}{r}
\;+\;\mathcal{O}(r^{-2}),
\end{equation}
with \(\overline{\mathbf{G}}_{\infty}\) transverse in the \(\mathbf{r}\)-direction
\cite{stratton1941}. Equation~\eqref{eq:ch2.sommerfeld-dyadic-asymptotic}
is the kernel form of Eq.~\eqref{eq:ch2.sommerfeld-field-condition} under
Eq.~\eqref{eq:ch2.electric-dyadic-convolution}: impressed currents that satisfy
Eq.~\eqref{eq:ch1.source-admissibility} generate fields whose far zone is built from
outgoing dyadic radiation \cite{stratton1941}. The leading \(r^{-1}\) term in
Eq.~\eqref{eq:ch2.free-space-dyadic-green} is the explicit vacuum instance of
Eq.~\eqref{eq:ch2.sommerfeld-dyadic-asymptotic} \cite{jackson1999}. Subsection~\ref{sec:ch242}
promotes Eq.~\eqref{eq:ch2.sommerfeld-dyadic-asymptotic} to an operator selector on
\(\mathcal{S}_\omega\); Subsection~\ref{sec:ch243} uses it in uniqueness energy
integrals \cite{harrington2001}.
Collecting the asymptotic and impedance conditions above yields the admissible function class
on which exterior uniqueness arguments and radiation export are defined. A field pair belongs
to this class only if it satisfies the vector Sommerfeld condition and the transverse far-zone
template simultaneously; incoming or standing-wave completions are excluded by construction
\cite{harrington2001}. Define the Sommerfeld outgoing subspace on an exterior shell
\(\Gamma_R=\{r=R\}\):
\begin{equation}
\label{eq:ch2.sommerfeld-outgoing-class}
\mathcal{F}_{\mathrm{SF}}
=
\left\{
(\E,\H)\in\HCurl(\Omega_{\mathrm{ext}}):
\text{Eqs.~\eqref{eq:ch2.sommerfeld-field-condition}--\eqref{eq:ch2.far-zone-transverse-radiation} hold}
\right\}.
\end{equation}
In Eq.~\eqref{eq:ch2.sommerfeld-outgoing-class},
\(\mathcal{F}_{\mathrm{SF}}\subset\HCurl(\Omega_{\mathrm{ext}})\) is the Sommerfeld
outgoing field class on exterior \(\Omega_{\mathrm{ext}}\): pairs \((\E,\H)\) whose
far-zone asymptotics satisfy outgoing phase and transverse radiation templates;
incoming and standing-wave completions are excluded by membership in this set.
Equation~\eqref{eq:ch2.sommerfeld-outgoing-class} embeds Eq.~\eqref{eq:ch1.radiation-trace-condition}
into the Chapter~\ref{ch:02} function-space setting of
Eq.~\eqref{eq:ch2.admissible-field-space} \cite{stratton1941}. The physical
solution operator \(\mathcal{S}_\omega\) maps \(\mathcal{J}_\Omega\) into
\(\mathcal{F}_{\mathrm{SF}}\) when \(\mathcal{B}_\Gamma\) includes outgoing radiation condition
\cite{harrington2001}. Incoming solutions lie outside
\(\mathcal{F}_{\mathrm{SF}}\) even if they satisfy volume Maxwell rows; that exclusion is the Sommerfeld radiation condition \cite{stratton1941}.
Finite computations cannot place a boundary at infinity; the Silver--M\"uller trace below
approximates the Sommerfeld impedance relation on a truncation sphere large enough that
spurious reflections are small compared with exported power \cite{harrington2001}. On a finite
truncation sphere \(\Gamma_R\), Eq.~\eqref{eq:ch2.sommerfeld-impedance-ratio}
motivates the first-order absorbing condition
\begin{equation}
\label{eq:ch2.silver-muller-trace-approx}
\hat{\mathbf{n}}\times\H
\;-\;
\eta_0^{-1}(\hat{\mathbf{n}}\times\E)\times\hat{\mathbf{n}}
=
\mathbf{0}\quad\text{on }\Gamma_R,
\end{equation}
which approximates Eq.~\eqref{eq:ch1.radiation-trace-condition} at finite \(R\)
\cite{stratton1941}. Equation~\eqref{eq:ch2.silver-muller-trace-approx}
is not identical to Sommerfeld completion at infinity: spurious reflections from
\(\Gamma_R\) contaminate near fields unless \(R\) is large compared with source
extent and wavelength \cite{harrington2001}. Section~\ref{sec:ch24} treats
Eq.~\eqref{eq:ch2.silver-muller-trace-approx} as a practical surrogate; the
continuous target remains Eq.~\eqref{eq:ch2.sommerfeld-field-condition}
\cite{stratton1941}. Perfectly matched layers and radiation boundary integrals are
discrete refinements of the same outgoing-trace logic \cite{harrington2001}.
Once Sommerfeld closure is in place, the radiation export map $\mathcal{R}_\omega$ and projector
$\Pi_{\mathrm{rad}}$ act on the same outgoing class the range attributes to measured patterns:
$\mathcal{R}_\omega[\Jimp]=\Pi_{\mathrm{rad}}\,\mathcal{S}_\omega[\Jimp]$ with full field pair
$\mathcal{S}_\omega[\Jimp]$ carrying reactive content that $\Pi_{\mathrm{rad}}$ removes before
far-zone metering \cite{stratton1941}. The outgoing selector \(\Pi_{\mathrm{rad}}\) of
Eq.~\eqref{eq:ch2.radiation-operator-def} is idempotent on fields satisfying
Eq.~\eqref{eq:ch2.sommerfeld-outgoing-class} once uniqueness is established in
Subsection~\ref{sec:ch243} \cite{stratton1941}. Prior to that proof, Sommerfeld
closure defines the domain on which \(\Pi_{\mathrm{rad}}\) acts as projection rather than heuristic windowing \cite{harrington2001}. Composition with
Eq.~\eqref{eq:ch2.transport-qualified-radiation} yields
\begin{equation}
\label{eq:ch2.sommerfeld-radiation-composition}
\mathcal{R}_\omega[\Jimp]
=
\Pi_{\mathrm{rad}}\,\mathcal{S}_\omega[\Jimp],
\qquad
\mathcal{S}_\omega[\Jimp]\in\mathcal{F}_{\mathrm{SF}},
\end{equation}
with transport qualification from Section~\ref{sec:ch01.2.4} applied after
Eq.~\eqref{eq:ch2.sommerfeld-radiation-composition} \cite{stratton1941}. Equation
\eqref{eq:ch2.sommerfeld-radiation-composition} expresses
Eq.~\eqref{eq:ch2.radiation-time-frequency-equiv} on the Sommerfeld domain without
re-proving causality from Subsection~\ref{sec:ch223} \cite{harrington2001}. Consequence for
range practice: pattern meters and gain reports must compose with $\mathcal{R}_\omega$, not with
raw $\mathcal{S}_\omega[\Jimp]$, whenever reactive near fields are non-negligible on the
measurement shell \cite{stratton1941}. Subsection~\ref{sec:ch241} therefore states
what the compact range already assumes before any far-zone integral is evaluated in
Section~\ref{sec:ch25} \cite{stratton1941}.
Equations~\eqref{eq:ch2.sommerfeld-scalar-asymptotic}--\eqref{eq:ch2.sommerfeld-outgoing-class}
require linear vacuum far zones, fixed \(\omega>0\), and compact \(\operatorname{supp}\Jimp\).
They fail for anisotropic or magneto-electric exterior media without modified
asymptotics, for lossless closed cavities where Eq.~\eqref{eq:ch1.outer-boundary-classes}
uses PEC rather than outgoing radiation condition, and for time-harmonic fields with net incoming
power at infinity \cite{stratton1941}. Loss shifts
Eq.~\eqref{eq:ch2.sommerfeld-scalar-asymptotic} to complex \(k_0\) but preserves
outgoing \emph{selection} as distinct from incoming \cite{stratton1941}.
As a worked example, replace Eq.~\eqref{eq:ch2.sommerfeld-field-condition} by
Eq.~\eqref{eq:ch2.silver-muller-trace-approx} on \(\Gamma_R\) with \(R=5\lambda_0\).
Fields inside \(\Omega_s\) differ from the exact Sommerfeld solution by
\(\mathcal{O}((k_0 R)^{-1})\) in typical radiative components while reactive near
content may differ by order unity \cite{stratton1941}. Declaring a
transport channel from \(\Gamma_R\) probes without
Eq.~\eqref{eq:ch2.regime-operator-screen} conflates truncation error with angular
transport \cite{stratton1941}.
